%% ma: L10-B14-0003
%% lop: 10
%% chuong: 5
%% bai: 14
%% dang: 10.14.2
%% loai: TL
%% muc: TH
%% dap_an: "-"
%% nguon: "chuyen-de-cac-so-dac-trung-cua-mau-so-lieu-khong-ghep-nhom-toan-10-chuong-trinh-moi.pdf (GV. Phan Nhật Linh) (P04), Bài 03 Dạng 1 Bài tập tự luận 3"
%% kien_thuc: "Lớp 10 Bài 13 (tứ phân vị), Bài 14 (khoảng tứ phân vị)"
%% trang_thai: chinh-thuc
%% ghi_chu: "Tự duyệt theo yêu cầu giáo viên 22/10/2026."
\begin{ex}
Điểm bài thi khảo sát môn Toán học kì 1 của tổ 1 lớp 10A như sau:
\begin{center}
\begin{tabular}{|c|c|c|c|c|c|c|c|c|c|c|c|}\hline
Tên & An & Nam & Hòa & Trung & Thủy & Tiến & Nguyên & Hồng & Lan & Thúy & Hà\\\hline
Điểm & $8$ & $7$ & $9$ & $8$ & $5$ & $7$ & $8$ & $10$ & $9$ & $8$ & $9$\\\hline
\end{tabular}
\end{center}
Tìm khoảng biến thiên, khoảng tứ phân vị của mẫu số liệu.
\loigiai{Sắp xếp: $5$; $7$; $7$; $8$; $8$; $8$; $8$; $9$; $9$; $9$; $10$. $R=10-5=5$.

$n=11$: $Q_2=8$; nửa dưới $5$; $7$; $7$; $8$; $8$ nên $Q_1=7$; nửa trên $8$; $9$; $9$; $9$; $10$ nên $Q_3=9$.

$\Delta_Q=9-7=2$.}
\end{ex}
